\section{案例 3：研究工作中的写作回避}

\subsection*{情境}
一名研究生能够读论文、讨论想法，但却反复拖延写作初稿。

\subsection*{框架映射}
\begin{itemize}[nosep]
  \item \textbf{当前状态：} 浏览、批注与讨论，但没有可见输出。
  \item \textbf{目标状态：} 产出粗糙但明确写下来的主张。
  \item \textbf{主导变量：} 伪工作状态里存在中等 $\Reward$，由于歧义导致极高的 $\Switch$，以及朝向低暴露任务的自我保护性环境漂移。
  \item \textbf{错过的决策窗口：} 写作会话的第一分钟，那时原本还可以把草稿狭义地定义出来。
\end{itemize}

\subsection*{机制解释}
这不是一个简单的“不活动”案例，而是一个相邻吸引子案例。学生其实一直在活动，但活动发生在另一个盆地里；那个盆地既能提供智识上的正当感，又不会迫使他把结构外显出来。涌现语言在这里重要，是因为反复在阅读或讨论里获得成功，会形成一种路径依赖的身份回路：系统越来越擅长草稿之前的机制，因此转入可见产出的成本也越来越高。吸引子语言在这里有帮助，因为问题不是完全不动，而是在状态空间的错误区域内部移动。

\subsection*{证据车道纪律}
这里的机制主张属于\textbf{中等}。吸引子与亚稳态相关文献支持把盆地、持续性和重构语言当作一种形式支架来使用。但这些文献并不能支持一个强经验主张，即我们已经可以直接从脑层级测量中读取“写作回避”。因此，这个案例仍然是一个遵守纪律的行为模型，它受系统概念支持，但并不是一个已经完成的实验室结论。

\subsection*{操作性修复}
\begin{enumerate}[nosep]
  \item 把目标改写成一个丑陋但可见的最小单元：五条要点主张、一段段落骨架，或一条图注。
  \item 以产出性工件而不是阅读性工件来开启会话。
  \item 让上一次会话遗留下来的一个未解决问题保持可见。
  \item 把评价与生成拆开，使句子产出在进入阶段不必同时承担完整的质量审判。
\end{enumerate}

\subsection*{迁移教训}
有些失败并不是因为系统完全瘫住，而是因为它在相邻状态里变得过于熟练。面对这种情况，干预需要降低目标歧义，并暴露出最小的真实输出单元，而不是仅仅延长停留在目标邻域附近的时间。